\documentclass[12pt,a4paper]{article}
\usepackage[margin=1in]{geometry}
\usepackage{amsmath,amssymb,amsthm}
\usepackage[hidelinks]{hyperref}
\usepackage{booktabs}
\usepackage{array}
\newcolumntype{P}[1]{>{\raggedright\arraybackslash}p{#1}}

\newtheorem{definition}{Definition}[section]
\newtheorem{remark}[definition]{Remark}

\title{Stable Hulls as Imprint Integration:\\
An Interpretive Bridge from Filtered Traces to Self-Reading World-Models}

\author{
  Sidong Liu, PhD \\
  iBioStratix Ltd \\
  \texttt{sidongliu@hotmail.com}
}

\date{July 2026}

\begin{document}
\emergencystretch=2em
\raggedbottom
\hbadness=5000
\vbadness=5000

\maketitle

\begin{abstract}
This paper develops a programme-level bridge between the imprint vocabulary of \emph{The World as Imprint} and the accessible-world machinery distributed across Foundation~II, Foundation~III, and \emph{The Brain as a Plastic Observable-Algebra Network}. Its claim is conditional and architectural. The decoder-relative imprint may be small or null, so neither carrying nor stable-hull formation is guaranteed. When readable traces are present, the paper proposes stable hulls as integration regimes through which selected traces may become accessible, retainable, reportable, and eventually available for organized self-reading. It does not identify stable hulls with consciousness or provide sufficient conditions for first-person subjectivity.

The central thesis is likewise conditional: a finite observer does not encounter a full candidate imprint-field, but only those readable traces, if any, that enter an accessible algebra, stabilize under repeated uptake, and become integrated into a self-world map. Stable hulls are proposed as a middle regime between candidate imprint-bearing structure and uneven self-reading. They describe how readable traces may become a usable world before becoming a reflexively read self. The contribution is internal architectural clarification for an existing multi-paper programme, not a new result in cognitive science or neuroscience.

\medskip

\noindent\textbf{Keywords}: stable hull, imprint integration, accessible algebra, self-model, memory, self-reading, world-model
\end{abstract}

\section{Introduction: From Imprint-Field to Accessible World}

\emph{The World as Imprint} argued that the emergent world can be read as an imprint-field: a domain of structures insofar as they carry readable marks of the generative constraints that shape them \cite{WorldImprint}. That heuristic explicitly allowed the decoder-relative imprint to be small or null; carrying is not guaranteed. But the same paper also insisted that carrying is not the same as reflexive recognition. Not every structure that bears a readable mark thereby counts as a mind. What remained underdeveloped was the middle regime between these two claims. How do readable traces become an accessible world for a finite observer before they become the object of organized self-reading?

This paper develops that middle regime through a distinct cognitive
object, the integration workspace \(\mathcal W_{\mathrm{int}}\). It is
not a mysterious substance or a synonym for consciousness. It is a
programme-level integration regime in which filtered traces are proposed
to become jointly available for memory, report, coordinated uptake, and
self-location.

The bridge matters because without it the path from imprint-carrying to self-reading remains too abrupt. One would then have only two levels: candidate imprint-bearing structure on one side and mind on the other. But the actual architecture of finite observers is richer. Before a system can read itself as imprint-bearing, it must inhabit a world that is stable enough to support retained salience, cross-temporal coordination, and self-world differentiation. That world is not the whole imprint-field. It is a stabilized accessible sector carved out of it.

Earlier papers use similar hull language for different typed objects.
Foundation~I has a source-algebraic intersection object,
Foundation~II has an external shared accessible sector, and the Brain
Anchor discusses an internal cognitive workspace
\cite{FoundationI,FoundationII,BrainAnchor}. The present paper concerns
only the third type, denoted \(\mathcal W_{\mathrm{int}}\). It does not
identify or derive it from either Foundation object.

\begin{quote}
filtered imprints $\to$ accessible organization $\to$
\(\mathcal W_{\mathrm{int}}\) $\to$ retained memory $\to$ self-world model.
\end{quote}

Several guardrails are necessary at the outset. First,
\(\mathcal W_{\mathrm{int}}\) is proposed as a precondition for organized
self-reading, not a sufficient condition for consciousness. Second,
memory integration is not personal survival or identity transfer. Third,
this paper does not determine which animals or machines are conscious.
Fourth, the self-reading threshold is treated elsewhere
\cite{WorldImprint}. The present paper is an interpretive companion, not
an independent contribution to self-model theory or neuroscience.

The paper proceeds accordingly. Section~2 separates the three hull
namespaces. Section~3 introduces filtered imprints and accessible
organization. Section~4 interprets \(\mathcal W_{\mathrm{int}}\) as a
cognitive integration regime. Later sections treat memory, self-world
modeling, interfaces to other papers, and limitations.

\section{Hull Namespace and Type Registry}

The phrase ``stable hull'' has been used for three objects with different
domains and types:

\begin{description}
\item[\(\mathcal H_{\mathrm{st}}\)] the Foundation~I
source-algebraic stable hull, defined by an intersection construction in
the source/semigroup setting \cite{FoundationI};
\item[\(\mathcal H_{12}\)] the Foundation~II external shared accessible
sector for two independently specified source chains
\cite{FoundationII};
\item[\(\mathcal W_{\mathrm{int}}\)] the internal cognitive integration
workspace used in this bridge and associated with the Brain Anchor
programme \cite{BrainAnchor}.
\end{description}

No equality, inclusion, or canonical map among these three objects is
established. In particular, \(\mathcal H_{12}\) is not an internal brain
workspace, and \(\mathcal W_{\mathrm{int}}\) is not the Foundation~I
intersection \(\mathcal H_{\mathrm{st}}\).

\begin{definition}[Integration workspace]
The \emph{integration workspace} \(\mathcal W_{\mathrm{int}}\) is a
programme-level cognitive regime in which a family of filtered traces is
sufficiently persistent, mutually coherent, and recurrently available
to support world-level uptake, retention, and coordinated redescription.
\end{definition}

Three phrases in the definition matter.

First, ``cognitive regime'' means that \(\mathcal W_{\mathrm{int}}\) is
typed relative to internal uptake by a finite system, unlike the external
agreement object \(\mathcal H_{12}\). Second, ``mutually coherent'' means
that retained traces support a structured world-picture. Third,
``recurrently available'' means that traces can be compared, updated,
and reused.

This definition is an architectural proposal, not a consequence of
Foundation~II. The imprint heuristic also does not guarantee that traces
exist everywhere: relative to a source family and decoder class the
imprint value may be zero. The question here is conditional---when
readable traces exist, which become organized enough to form a usable
world-model?

\section{Filtered Imprints and Accessible Algebra}

The first bridge paper argued that coarse-graining can be read as imprint filtering: not every lower-level feature survives scale change, and those that do survive constitute scale-stable residues. But even after such filtering, the resulting imprint-field is still too large to be inhabited in full by any finite observer. A second selection is therefore necessary. This second selection is not primarily a matter of scale; it is a matter of access.

\begin{definition}[Accessible imprint]
An \emph{accessible imprint} is a filtered trace that lies within the observer-relative regime of uptake, retention, comparison, or action-guidance.
\end{definition}

This definition is deliberately observer-relative without becoming subjectivist. The point is not that accessibility is arbitrary or purely private. The point is that finite systems encounter only a structured subset of all available marks. An accessible algebra is one way to formalize this subset in the broader programme \cite{FoundationI,FoundationII}. For present purposes, the essential idea is simpler: the world of a finite observer is formed by those traces that can enter stable uptake.

This gives a two-step selection picture:

\begin{enumerate}
\item coarse-graining selects which lower-level features survive as scale-stable imprints;
\item accessibility selects which of those imprints can enter an observer's usable world.
\end{enumerate}

The proposed workspace \(\mathcal W_{\mathrm{int}}\) lives at the
second step. It does not explain why a trace survives scale change. It
names the further cognitive integration problem for a readable trace.

This is also where the bridge differs from standard cognitive-theory usage. It does not claim to replace global workspace, predictive processing, or enactive models. It claims only to identify an architectural slot that these literatures often presuppose in different ways: the transition from mere accessible traces to a stable, reusable world-organization.

\section{The Integration Workspace \texorpdfstring{\(\mathcal W_{\mathrm{int}}\)}{W-int}}

It is tempting to think that once traces are accessible, a world is already present. That is too quick. Accessibility alone may still be fragmented, transient, or mutually inconsistent. What is needed for worldhood is integration.

\subsection{Why accessibility is not enough}

A finite system can register signals without organizing them into a
stable world. It may encounter accessible imprints but fail to build an
integration workspace \(\mathcal W_{\mathrm{int}}\). The distinction
marks the difference between signal-response and the proposed
architecture of inhabitable worldhood.

\subsection{Integration conditions}

Without pretending to state a theorem, the following conditions capture
what the \(\mathcal W_{\mathrm{int}}\) bridge is trying to mark:

\begin{center}
\begin{tabular}{P{0.28\linewidth}P{0.6\linewidth}}
\toprule
\textbf{Condition} & \textbf{Role in \(\mathcal W_{\mathrm{int}}\) integration} \\
\midrule
persistence & the same or comparable traces remain available across time \\
coherence & retained traces support a jointly navigable picture rather than disconnected shards \\
selective retention & some traces are stabilized while others are dropped or backgrounded \\
cross-context usability & traces can guide report, action, or re-identification across situations \\
self-world differentiation & the regime can support a distinction between world-features and self-locating features \\
\bottomrule
\end{tabular}
\end{center}

None of these conditions is by itself equivalent to consciousness. Together they specify something weaker but still crucial: a world-forming regime in which accessible traces become jointly organized. These conditions also resonate with global-availability accounts of conscious access \cite{Baars1988,DehaeneChangeux2011}. In that sense, the table is best read as an interface map rather than as a rival empirical theory.

\subsection{Integration workspace and reportable worldhood}

The integration workspace is proposed as a precondition for reportable
worldhood. A system may have filtered inputs without such a workspace;
in that case the bridge supplies no stable map for comparison or
re-entry. Even when \(\mathcal W_{\mathrm{int}}\) is present, reflexive
self-reading does not follow.

\section{Memory as Retained Imprint Integration}

Memory illustrates why a distinct cognitive integration interface is
useful. It is selective retention for later integration, not an instance
of the external agreement sector \(\mathcal H_{12}\).

\subsection{Retention is selective}

No finite system copies a full source description. Memory operates by
selective retention: some traces stabilize, others decay or are
overwritten, and some remain only as biasing conditions. What survives
may be re-entered into \(\mathcal W_{\mathrm{int}}\).

\subsection{Memory is integrative, not archival}

The metaphor of archive is misleading. Memory is not a warehouse of preserved states. It is better understood as retained imprint integration. A memory trace matters not because it exists in isolation, but because it can be recruited back into a larger world-model. It contributes to expectations, error-correction, self-location, and salience structure. Within predictive-processing accounts, retained traces function as priors shaping anticipation and prediction-error correction \cite{Clark2013}.

The Brain Anchor paper motivates an internal workspace interpretation
through plasticity and recurrent world-modeling \cite{BrainAnchor}. The
present paper does not verify those claims or import them as a theorem;
their consistency with the namespaced \(\mathcal W_{\mathrm{int}}\) is a
separate interface question. Plastic stabilization is used here only as
a candidate illustration of retained world-structure.

\subsection{No claim about identity transfer}

The memory discussion also demands a guardrail already familiar from \emph{The World as Imprint}. Retention of traces does not establish persistence of the same first-person subject elsewhere \cite{WorldImprint}. The fact that effects propagate, habits persist, or structures re-enter later organization should not be confused with personal identity transfer. In the present framework, memory integration concerns world-organization, not metaphysical survival.

\section{The Self-World Model as Organized Imprint Map}

Once an integration workspace \(\mathcal W_{\mathrm{int}}\) is
stipulated, a further transition becomes possible: construction of a
self-world model. This is an architectural ordering, not a theorem of
cognitive necessity.

\subsection{Map, not inner substance}

The phrase ``self-world model'' risks misunderstanding if treated too quickly. It does not refer to a ghostly entity observing the world from behind the scenes. It names an organized map in which world-features, bodily regularities, action possibilities, and self-locating coordinates become mutually articulated. In imprint language, it is a way of organizing integrated traces into a differentiated picture. This formulation aligns with the self-model theory of subjectivity \cite{Metzinger2003}, while still treating the present paper's claim as an architectural redescription rather than a reconstruction of phenomenal transparency.

\subsection{Why the self-model comes after the workspace in this bridge}

Within this proposed bridge, a self-model is placed after a sufficiently
stable world-map. Accordingly \(\mathcal W_{\mathrm{int}}\) is a
stipulated architectural precondition for the self-model, not a result
inherited from \(\mathcal H_{\mathrm{st}}\) or \(\mathcal H_{12}\). The
bridge stops before the self-reading threshold.

\subsection{Organized imprint map}

The self-world model can therefore be described in a sentence:

\begin{quote}
A self-world model is an organized imprint map in which retained traces are structured into a navigable distinction between self-locating, world-locating, and action-guiding organization.
\end{quote}

This formulation does not yet say that anything ``feels like something.'' It says only that a finite regime has reached the point at which world-organization and self-location are jointly stabilized. Whether the integration is best described in narrative terms or in transparent representational terms remains open \cite{Dennett1992,Metzinger2003}.

\section{Relation to Existing Papers}

The present bridge is useful only if its place in the larger body of work is explicit.

\subsection{Relation to Foundation~II}

Foundation~II studies the external shared accessible sector
\(\mathcal H_{12}\) across independently specified source-side chains
\cite{FoundationII}. That object is not the internal cognitive workspace
\(\mathcal W_{\mathrm{int}}\). The former proposal to import
Foundation~II as a formal derivation or de-tautologization of the present
workspace is retracted. At most, external agreement motivates an analogy;
no typed bridge has been constructed.

\subsection{Relation to Foundation~III}

Foundation~III adds conditional feedback structure at the
source-manifest interface \cite{FoundationIII}. The present paper imports
no feedback theorem. It uses historically shaped organization only as a
programme-level analogy for retention and re-entry in
\(\mathcal W_{\mathrm{int}}\).

\subsection{Relation to Brain Anchor}

The Brain Anchor paper offers a candidate biological realization of an
internal integration workspace \cite{BrainAnchor}. The present paper
uses the shared symbol \(\mathcal W_{\mathrm{int}}\) for that cognitive
type but does not verify the published paper's internal consistency or
identify it with either Foundation hull. That verification is a separate
Gate~2 interface task.

\subsection{Relation to \emph{The World as Imprint}}

Finally, the relation to \emph{The World as Imprint} is direct. That paper distinguished decoder-relative carrying, which may be small or null, from uneven recognition \cite{WorldImprint}. The present paper fills the middle gap. It explains how carrying becomes accessible worldhood before it becomes reflexive recognition.

\section{Output to the Self-Reading Threshold}

The bridge can now be stated in its cleanest form:

\begin{quote}
Within this bridge, \(\mathcal W_{\mathrm{int}}\) is a proposed
precondition for organized self-reading, not a sufficient condition for
consciousness.
\end{quote}

This sentence is the paper's most important guardrail. It says what the paper contributes and where it must stop.

What the paper contributes is a clearer picture of the middle architecture:

\begin{enumerate}
\item filtered imprints survive scale change;
\item some of them enter accessible uptake;
\item some of those are proposed to enter \(\mathcal W_{\mathrm{int}}\);
\item that workspace supports the architectural description of retained
memory and a self-world model;
\item only then does the question of self-reading threshold arise.
\end{enumerate}

What the paper does not contribute is a theory of why first-person experience accompanies such organization. Nor does it provide a criterion by which all borderline biological or artificial cases can be settled. Those are later or different questions.

This is why the paper's stopping point matters. It stops at the integration of filtered imprints into an accessible self-world model; the self-reading threshold itself is treated elsewhere \cite{WorldImprint}. That is not a deficiency. It is the correct division of labor.

\section{Limitations}

Several limitations should remain explicit.

First, the paper is programme-level. It does not construct
\(\mathcal W_{\mathrm{int}}\) from arbitrary imprint-fields.

Second, the terminology of hull, access, retention, and model remains partly architectural rather than fully operationalized. Further formal work would need metrics for integration depth, accessibility bandwidth, and retention stability.

Third, the bridge remains neutral on substrate. It does not assume that
only brains can realize \(\mathcal W_{\mathrm{int}}\), but it certifies
no biological or non-biological system as conscious.

Fourth, the paper does not decide whether systems realizing the proposed
workspace cross the self-reading threshold. The two questions remain
separate.

Fifth, the paper should not be read as offering a new empirical discriminator among existing self-model theories. Its role is interpretive and organizational: it clarifies how filtered imprints, accessible organization, and self-world mapping fit together inside the larger programme.

Despite those limitations, the namespace makes the proposed transition
more precise without claiming formal inheritance:

\begin{quote}
readable imprints $\to$ accessible world $\to$
\(\mathcal W_{\mathrm{int}}\) $\to$ self-world model $\to$
self-reading threshold.
\end{quote}

This remains an interface architecture, not a theorem that the workspace
exists or produces a mind.

\begin{thebibliography}{99}

\bibitem{WorldImprint}
S.~Liu, \emph{The World as Imprint: From Trace-Carrying Structure to Self-Reading Mind}, Zenodo (2026), DOI: 10.5281/zenodo.20122718.

\bibitem{FoundationI}
S.~Liu, \emph{Foundation I: From Source to Manifestation --- A Conditional Same-Root Program for Geometry and Modular Flow}, Zenodo (2026), DOI: 10.5281/zenodo.19627786.

\bibitem{FoundationII}
S.~Liu, \emph{Foundation II: Shared Manifestation --- A Conditional Two-Source Program for Shared Accessible Sectors}, Zenodo (2026), DOI: 10.5281/zenodo.19641372.

\bibitem{FoundationIII}
S.~Liu, \emph{Foundation III: The Perfuming Loop --- A Conditional One-Step Source-Level Feedback Program}, Zenodo (2026), DOI: 10.5281/zenodo.19642189.

\bibitem{BrainAnchor}
S.~Liu, \emph{The Brain as a Plastic Observable-Algebra Network: The Conditional Fixed-Point Workspace $\mathcal{W}_{\mathrm{int}}$}, Zenodo (2026), DOI: 10.5281/zenodo.20044329.

\bibitem{Baars1988}
B.~J.~Baars,
\emph{A Cognitive Theory of Consciousness},
Cambridge University Press, Cambridge (1988).

\bibitem{DehaeneChangeux2011}
S.~Dehaene and J.-P.~Changeux,
``Experimental and theoretical approaches to conscious processing,''
\emph{Neuron} \textbf{70} (2011), 200--227,
DOI: 10.1016/j.neuron.2011.03.018.

\bibitem{Metzinger2003}
T.~Metzinger,
\emph{Being No One: The Self-Model Theory of Subjectivity},
MIT Press (2003).

\bibitem{Dennett1992}
D.~C.~Dennett,
\emph{The self as a center of narrative gravity},
in F.~Kessel, P.~Cole, and D.~Johnson (eds.),
\emph{Self and Consciousness: Multiple Perspectives},
Lawrence Erlbaum, 103--115 (1992).

\bibitem{Clark2013}
A.~Clark,
\emph{Whatever next? Predictive brains, situated agents, and the future of cognitive science},
Behavioral and Brain Sciences \textbf{36}, 181--204 (2013).

\bibitem{Thompson2007}
E.~Thompson,
\emph{Mind in Life: Biology, Phenomenology, and the Sciences of Mind},
Harvard University Press (2007).

\end{thebibliography}

\end{document}
